\documentclass{article}

\usepackage{amsmath,amssymb}
\usepackage{xurl}
\usepackage[hidelinks]{hyperref}
\setlength{\emergencystretch}{2em}

% Shared commands for notes in docs/paper-gaps/.

\DeclareUnicodeCharacter{2080}{\ifmmode{}_{0}\else\textsubscript{0}\fi}
\DeclareUnicodeCharacter{2081}{\ifmmode{}_{1}\else\textsubscript{1}\fi}
\DeclareUnicodeCharacter{2082}{\ifmmode{}_{2}\else\textsubscript{2}\fi}
\DeclareUnicodeCharacter{2099}{\ifmmode{}_{n}\else\textsubscript{n}\fi}

\newcommand{\Lean}{\textsc{Lean}}
\ExplSyntaxOn
\NewDocumentCommand{\leanid}{m}{
  \ifmmode
    \textsf{\detokenize{#1}}
  \else
    \group_begin:
    \urlstyle{sf}
    \tl_set:Nn \l_tmpa_tl {#1}
    \tl_replace_all:Nnn \l_tmpa_tl {\_} {_}
    \exp_args:NV \path \l_tmpa_tl
    \group_end:
  \fi
}
\ExplSyntaxOff
\providecommand{\tr}{\operatorname{tr}}
\newcommand{\tnleanIssueBase}{https://github.com/LionSR/TNLean/issues/}
\newcommand{\tnleanPrBase}{https://github.com/LionSR/TNLean/pull/}
\newcommand{\tnleanDocBase}{https://sirui-lu.com/TNLean/docs/}
\newcommand{\ghissue}[1]{\href{\tnleanIssueBase#1}{GitHub issue \##1}}
\newcommand{\ghpr}[1]{\href{\tnleanPrBase#1}{PR \##1}}
\newcommand{\doclink}[1]{\href{\tnleanDocBase#1.html}{\path{#1}}}

% Dirac notation and the identity operator, as in blueprint/src/macros/common.tex.
% \providecommand keeps note-local definitions authoritative.
\usepackage{dsfont}
\providecommand{\ket}[1]{|#1\rangle}
\providecommand{\bra}[1]{\langle #1|}
\providecommand{\braket}[2]{\langle #1 | #2 \rangle}
\providecommand{\ketbra}[2]{|#1\rangle\!\langle #2|}
\providecommand{\Id}{\mathds{1}}
\providecommand{\one}{\mathds{1}}
\providecommand{\C}{\mathbb{C}}
\providecommand{\MN}[1]{M_{#1}(\C)}
\providecommand{\MD}{M_D(\C)}

% External referencing for paper sources.  The build helper writes sanitized
% auxiliary files under build/paper-aux/ and excludes duplicated source labels.
\usepackage{xr}

% Import a generated paper auxiliary file with a caller-supplied label prefix.
% The candidate paths support builds from docs/paper-gaps/, the repository root,
% and build/paper-aux/ itself.
\newcommand{\importpaperaux}[2]{%
  \IfFileExists{../../build/paper-aux/#2.aux}{%
    \externaldocument[#1]{../../build/paper-aux/#2}%
  }{%
    \IfFileExists{build/paper-aux/#2.aux}{%
      \externaldocument[#1]{build/paper-aux/#2}%
    }{%
      \IfFileExists{#2.aux}{%
        \externaldocument[#1]{#2}%
      }{}%
    }%
  }%
}

% General external reference macros
\newcommand{\paperref}[2]{\ref{#1-#2}}
\newcommand{\papereqref}[2]{(\ref{#1-#2})}

% CPSV16 source (1606.00608 / MPDO-22-12-17-2.tex) external document and shortcuts
\importpaperaux{cpsv16-}{1606.00608/MPDO-22-12-17-2-tnlean}
\newcommand{\cpsvref}[1]{\paperref{cpsv16}{#1}}
\newcommand{\cpsveqref}[1]{\papereqref{cpsv16}{#1}}

% Verdict marker: \gapnote{<kind>}{<status>}. Kind is the policy
% classification (clarification, local-correction, scope-restriction,
% unfaithful, false-source, open-gap); status is open, wip (elimination
% underway), resolved, or historical. Machine-read by the site index;
% prints a verdict line.
\newcommand{\gapnote}[2]{\par\noindent\textbf{Verdict:} #1 (#2).\par}


\title{The Figure 11 Fusion Map Is a Coisometry onto the Active Sectors}
\author{}
\date{2026-07-20}

\begin{document}
\maketitle
\gapnote{local-correction}{resolved}

\section{Orientation of the fusion map}

The canonical-form convention of Ref.~\cite{Cirac2016MPDO_arXiv} retains only
nonzero blocks. Let $D_k$ denote their dimensions and let $D$ denote the
original dimension. The source states at lines~214--225 that
\begin{align}
    \sum_k D_k
        &\leq D.
    \label{eq:cpsv16_figure11_fusion_coisometry_dimension_bound}
\end{align}

Proposition~4.13 of Ref.~\cite{Cirac2016MPDO_arXiv} uses a matrix $U$ in the
orientation
\begin{align}
    U\widetilde M U^\dagger
        &= \bigoplus_\alpha \mu_\alpha\otimes M_\alpha.
    \label{eq:cpsv16_figure11_fusion_coisometry_vertical_form}
\end{align}
This identity appears at lines~945--951 and~1863--1870. Thus $U$ maps the
original space onto the retained direct sum. In standard terminology, $U$ is
a coisometry. If $P_{\mathrm{act}}$ denotes the projection onto the nonzero
canonical sectors, then
\begin{align}
    UU^\dagger
        &= I,
    \label{eq:cpsv16_figure11_fusion_coisometry_vertical_row_identity}\\
    U^\dagger U
        &= P_{\mathrm{act}}.
    \label{eq:cpsv16_figure11_fusion_coisometry_vertical_active_projection}
\end{align}
This orientation and the need for exact reconstruction are established in
\path{docs/paper-gaps/cpgsv17_vertical_isometry_zero_sector.tex}.

Theorem~4.14(iii) of Ref.~\cite{Cirac2016MPDO_arXiv} has the same orientation.
For each pair of basic normal tensor labels $\alpha,\beta$, define the product
tensor $P_{\alpha,\beta}=M_\alpha M_\beta$. The theorem states
\begin{align}
    U_{\alpha,\beta}P_{\alpha,\beta}^{ij}
        U_{\alpha,\beta}^\dagger
        &= D_{\alpha,\beta}^{ij},
    \label{eq:cpsv16_figure11_fusion_coisometry_fusion_forward}\\
    D_{\alpha,\beta}^{ij}
        &= \bigoplus_\gamma
            \chi_{\alpha,\beta,\gamma}\otimes M_\gamma^{ij}.
    \label{eq:cpsv16_figure11_fusion_coisometry_active_direct_sum}
\end{align}
These statements occur at lines~986--993. Appendix~C.4 constructs the
right-hand side from the active canonical blocks of the Figure~11 tensor
\cite[lines~2020--2029]{Cirac2016MPDO_arXiv}. The source-compatible
orthogonality relation is therefore
\begin{align}
    U_{\alpha,\beta}U_{\alpha,\beta}^\dagger
        &= I,
    \label{eq:cpsv16_figure11_fusion_coisometry_fusion_row_identity}
\end{align}
not $U_{\alpha,\beta}^\dagger U_{\alpha,\beta}=I$ on the entire product bond
space.

The forward identity alone does not ensure exact recovery of the product
tensor. The active decomposition must also satisfy
\begin{align}
    P_{\alpha,\beta}^{ij}
        &= U_{\alpha,\beta}^\dagger
            D_{\alpha,\beta}^{ij}U_{\alpha,\beta}.
    \label{eq:cpsv16_figure11_fusion_coisometry_fusion_reconstruction}
\end{align}
This identity states that the omitted common corner is zero. Together,
Eqs.~\ref{eq:cpsv16_figure11_fusion_coisometry_fusion_row_identity}
and~\ref{eq:cpsv16_figure11_fusion_coisometry_fusion_reconstruction} imply
both zipper identities and preserve every nonempty closed word.

\section{A zero product sector}

Let the physical index take the values zero and one. Consider two
one-dimensional classical sectors. The first has its only nonzero letter at
$(0,0)$, while the second has its only nonzero letter at $(1,1)$. Both are
nonzero normal tensors. Their product tensor nevertheless vanishes:
\begin{align}
    P_{0,1}^{ij}
        &= \sum_{q=0}^{1}M_0^{iq}\otimes M_1^{qj},
    \label{eq:cpsv16_figure11_fusion_coisometry_cross_product_sum}\\
    P_{0,1}^{ij}
        &= 0
    \label{eq:cpsv16_figure11_fusion_coisometry_zero_cross_product}
\end{align}
for every $i,j$, because the first factor requires $q=0$ and the second
requires $q=1$. This is the cross-sector product in the elementary classical
GHZ family.

For this pair, every retained multiplicity space may have dimension zero. The
unique map from the one-dimensional product bond space to the empty retained
space satisfies
Eqs.~\ref{eq:cpsv16_figure11_fusion_coisometry_fusion_row_identity},
\ref{eq:cpsv16_figure11_fusion_coisometry_fusion_forward},
and~\ref{eq:cpsv16_figure11_fusion_coisometry_fusion_reconstruction}. It
cannot satisfy
\begin{align}
    U^\dagger U
        &= I.
    \label{eq:cpsv16_figure11_fusion_coisometry_excluded_column_identity}
\end{align}
Requiring this identity on the one-dimensional product space would impose a
full-support hypothesis absent from Theorem~4.14 and Appendix~C.4 of
Ref.~\cite{Cirac2016MPDO_arXiv}.

\section{Formal realignment}

The structure \leanid{MPOTensor.BNTFusionIsometryFamily} assumes
$U_{\alpha,\beta}^\dagger U_{\alpha,\beta}=I$. It remains a valid conditional
full-support structure, but it is not the unrestricted conclusion of the
source theorem. It is retained only for results that assume full support
explicitly.

The source-facing construction tracked in \ghissue{4588} uses
\leanid{MPOTensor.BNTFusionCoisometryFamily}. Its data consist of:
\begin{enumerate}
    \item positive diagonal matrices
        $\chi_{\alpha,\beta,\gamma}$, including zero-dimensional multiplicity
        spaces;
    \item a coisometry $U_{\alpha,\beta}$ from each product bond space onto
        the retained nonzero direct sum;
    \item the forward identity in
        Eq.~\ref{eq:cpsv16_figure11_fusion_coisometry_fusion_forward};
    \item the reconstruction identity in
        Eq.~\ref{eq:cpsv16_figure11_fusion_coisometry_fusion_reconstruction}.
\end{enumerate}
The projector-closure canonical decomposition supplies these data directly.
The nonzero reducing corners give isometric inclusions into the product bond
space, and their adjoints form the required row coisometry. This construction
does not require the active projection to equal the identity.

The tensor-attached structure
\leanid{MPOTensor.BNTFusionTensorClause} now implements the same active-support
contract. It stores $UU^\dagger=I$, the forward fusion identity, and exact
reconstruction. The implication from the renormalization fixed-point maps is
proved by
\leanid{MPOTensor.HasBNTFusionTensorClause.of_isRFPViaTS}, completing the
realignment tracked in \ghissue{4225}. Conversion to the stronger structure
\leanid{MPOTensor.BNTFusionIsometryFamily} is allowed only under the additional
condition $U^\dagger U=I$, or an equivalent statement that the product tensor
has no discarded zero reducing corner.

\section{Verdict}

The phrase ``isometries $U_{\alpha,\beta}$'' in
Ref.~\cite{Cirac2016MPDO_arXiv} uses the retained-row orientation already used
in Proposition~4.13. Reading it as $U^\dagger U=I$ strengthens the theorem and
excludes valid zero product sectors. The faithful Figure~11 conclusion is a
coisometry onto the active direct sum together with exact reconstruction. The
source-facing tensor-attached fusion clause and its renormalization-fixed-point
construction now implement exactly this conclusion.

\bibliographystyle{alpha}
\bibliography{references}

\end{document}
